\chapter{概率、信息与对偶工具}
\label{app:tools}

\section{事件数据处理}
对于分布$P,Q$和事件$E$，将和式分别限制到$E$与其补集，条件相对熵的非负性给出
\[
 \KL(P\Vert Q)\ge\kl(P(E),Q(E)).
\]
更一般的确定性分组与随机后处理，可以通过加入共同的随机核再条件分解得到。对二元相对熵求导，
\[
 \partial_p\kl(p,q)=\ln\frac{p(1-q)}{q(1-p)},\qquad
 \partial_q\kl(p,q)=\frac{q-p}{q(1-q)}.
\]
当$p>q$时，前者为正、后者为负。因此$p\ge1-\delta$、$q\le\delta$时最小值在边界取得。

\section{独立有界和的集中}
令$X\in[0,1]$，$f(\lambda)=\ln\EE e^{\lambda(X-\EE X)}$。有$f(0)=f'(0)=0$，二阶导为指数倾斜分布中的方差，至多$1/4$。两次积分得$f(\lambda)\le\lambda^2/8$。

对独立的$X_1,\ldots,X_s$，指数矩相乘。Markov不等式与参数最优化给出
\[
 \PP\!\left(\frac1s\sum_i(X_i-\EE X_i)\le-\gamma\right)
 \le e^{-2s\gamma^2}.
\]
这就是本书多数纠错使用的有界集中工具\citep{hoeffding1963}。共同偏差不会满足以真实目标为中心的重复正确性假设；必须先由观察模型确定多数在估计哪个变量。

\section{非负鞅的随时界}
设$M_t\ge0$为初值1的鞅。令$\tau$为首次$M_t\ge c$的时刻。对有限$n$，有$\EE M_{\tau\wedge n}=1$，而在$\tau\le n$时该变量至少$c$。因此$\PP(\tau\le n)\le1/c$。令$n$趋于无穷，
\[
 \PP(\sup_tM_t\ge c)\le1/c.
\]
对情境似然比，给定过去和所选实验，其下一乘因子的条件期望为$\sum_vK_a(v\mid\theta')=1$，故确为鞅。这证明通信停止规则所需的全部条件。非负超鞅的同类论证见\citet{howard2020}。

\section{条件Fano界}
令$Y$有$q$种可能值，估计由记录$W$生成，错误指示为$E$。给定$W$，当$E=0$时目标确定；当$E=1$时最多有$q-1$个可能值。因此
\[
 H_2(Y\mid W)\le h_2(p_e)+p_e\log_2(q-1).
\]
右侧在$0\le p_e\le1-1/q$上递增，故可用该区间内的错误上界$\delta$替换。令$W=(M,Z)$并减去条件熵，就得到正文的信息下界。独立解码随机性可以并入$W$，不会额外提供目标信息。

\section{有限覆盖规划的对偶}
把每个实验写成一个向量，其坐标是各冲突对的区别量；令$C$为这些向量的凸包。覆盖率是$\max_{x\in C}\min_jx_j$。对任意坐标概率$\lambda$，有$\min_jx_j\le\lambda^\top x$，这先给出原问题不大于对偶。

取$t>\kap$。紧凸集$C$与上正交区域$U_t=\{x:x_j\ge t\}$不相交。有限维凸集分离定理给出一个严格分离方向$\lambda$。由于$U_t$沿每个正坐标无界，分离方向必须非负且不全零。归一化$\sum_j\lambda_j=1$后，
\[
 \max_{x\in C}\lambda^\top x<t.
\]
让$t\downarrow\kap$并用概率单纯形的紧性，得到对偶不大于原问题。于是二者相等。该论证只使用有限维凸分离，没有给数值求解器的精度附加假设。
